%&pdftex

\def\proof{{\it Proof.} }
\font\msym=msbm10
\def\R{\hbox{\msym R}}
\def\N{\hbox{\msym N}}
%\def\R{\hbox{\bf R}}
%\def\N{\hbox{\bf N}}
\magnification=1200
\outer\def\art #1;#2;#3;#4;#5;#6;#7\par{\item {\bf [#1]} {\rm #2,{\it
#3},#4{\bf #5} (#6),#7}}

\centerline{\bf ASYMPTOTIC BEHAVIOUR OF ORTHOGONAL POLYNOMIALS}
\centerline{\bf RELATIVE TO MEASURES WITH MASS POINTS}

\bigskip

\centerline{Jos\'e J. Guadalupe
   \footnote{$^1$}{Dpto. de Matem\'aticas. Universidad de Zaragoza. 50009
             Zaragoza. Spain.},
Mario P\'erez $^1$,  Francisco J. Ruiz $^1$ and  Juan L. Varona
   \footnote{$^2$}{Dpto. de Matem\'atica Aplicada. Colegio Universitario de La
Rioja. 26001 Logro\~no. Spain.}}
\footnote{}{A.M.S. Classification: 42C05, 42C10.}
\footnote{}{The authors were partially supported by CAICYT, PB85--0338}
\bigskip

\footnote{}{AN UPDATED VERSION OF THIS MANUSCRIPT WERE PUBLISHED IN:
Mathematika {\bf 40} (1993), 331--344.}

{\narrower\narrower{\bf Abstract. }{\sl General expressions are found for the
orthonormal polynomials and the kernels relative to measures on the real line of
the form $d\mu+M\delta_c$, in terms of those of the measures $d\mu$ and
$(x-c)^2d\mu$. In particular, these relations allow to obtain several bounds for
the polynomials and kernels relative to a generalized Jacobi weight with a
finite number of mass points.}\smallskip}

\beginsection \S 0. Introduction.

Let $\mu$ be a positive measure on $\R$ with infinitely many points of increase
and such that all the moments
$$
   \int_{\R} x^nd\mu \qquad (n = 0,1,\dots)
$$
exist. Then, there exists a unique sequence $\{P_n\}_{n \geq 0}$ of orthogonal
polynomials
$$
   P_n(x) = k_n x^n + \dots, \qquad k_n>0
$$
such that
$$
   \int_{\R} P_n P_m d\mu = \cases{0,&if $n \not= m$; \cr 1, &if $n=m$.}
$$

As usual, $\{K_n(x,y)\}_{n \geq 0}$ denotes the sequence of kernels associated
to $\mu$, that is,
$$
   K_n(x,y) = \sum_{j=0}^n P_j(x) P_j(y).
$$

An interesting problem in the theory of orthogonal polynomials is that of
finding asymptotic estimates for $\{P_n\}$, their leading coefficients
$\{k_n\}$, the sequence $\{K_n(x,x)\}$ $(x \in {\rm supp }\mu)$, etc. (see, for
example, [8] and [4] for Jacobi polynomials, [2] and [6] for generalized Jacobi
polynomials, [1] and [5] for Laguerre and Hermite; general results can be
found in [7], [6], etc.).

Let $M$ be a positive constant and let $\delta_c$ denote a Dirac measure on a
point $c \in \R$, that is,
$$
   \int_{\R} f \delta_c = f(c)
$$
for every function $f$. Then, associated to the measure $d\nu = d\mu + M
\delta_c$ there exists a sequence $\{Q_n\}_{n \geq 0}$ of orthogonal
polynomials. We are interested in finding expressions which relate the sequences
$\{Q_n\}$ and $\{P_n\}$ in order to deduce estimates for $\{Q_n\}$ whenever
they are known for $\{P_n\}$. We will prove some relations which also involve
the polynomials orthonormal with respect to the measure $(x-c)^2d\mu(x)$. The
following notation will be used from now on:
$$
   d\mu^c = (x-c)^2d\mu(x);
$$

$\{P_n^c\}$ is the sequence of orthonormal polynomials relative to $d\mu^c$;
$$
   P_n^c(x) = k_n^c x^n + \dots, \qquad k_n^c > 0;
$$

$\{K_n^c(x,y)\}$ is the sequence of kernels relative to $d\mu^c$.

\beginsection \S 1. General results.

The following lemma gives a first relation among the sequences $\{K_n\}$,
$\{P_n\}$ and $\{P_n^c\}$:

\proclaim {Lemma 1}. With the above notation,
$$
   K_n(x,c) = {k_n \over {k_n^c}} P_n(c) P_n^c(x) - {k_{n-1}^c \over {k_{n+1}}}
   P_{n+1}(c) P_{n-1}^c (x) \qquad \forall n \geq 1.
$$

\proof We can put
$$
   K_n(x,c) = \sum_{j=0}^n \alpha_j P_j^c(x),
$$
with
$$
   \alpha_j = \int_{\R} K_n(x,c) P_j^c(x) (x-c)^2 d\mu(x).
$$
Therefore, we only need to show that

\centerline {a) $\alpha_n = {k_n \over k_n^c} P_n(c)$;\qquad\quad
             b) $\alpha_{n-1}=-{k_{n-1}^c \over k_{n+1}} P_{n+1}(c)$;\qquad\quad
             c) $\alpha_j=0$, $j=0,1,\dots,n-2$.}

Part a) can be obtained by a direct observation of the leading coefficients. As
to c), it is an easy consequence of a well-known property of the kernels $K_n$:
if $R_n$ is a polynomial of degree at most $n$, then
$$
   \int_{\R} K_n(x,c) R_n(x) d\mu(x) = R_n(c).
$$
In order to prove equation b), we use Christoffel-Darboux formula (see [8], for
example) and the orthonormality of $\{P_n\}$ with respect to $d\mu$:
$$ \eqalign{
   \alpha_{n-1} &= \int_{\R} K_n(x,c) P_{n-1}^c(x) (x-c)^2 d\mu(x) =
   \int_{\R} [K_n(x,c) (x-c)] [P_{n-1}^c(x) (x-c)] d\mu(x)
   \cr
   &= \int_{\R} {k_n \over k_{n+1}} [P_n(c)P_{n+1}(x) - P_{n+1}(c)P_n(x)]
   [P_{n-1}^c(x) (x-c)] d\mu(x)
   \cr
   &= -{k_n \over k_{n+1}}P_{n+1}(c) \int_{\R} P_n(x)[P_{n-1}^c(x) (x-c)] d\mu(x)
   = - {k_{n-1}^c \over k_{n+1}} P_{n+1}(c),
   \cr}
$$
and the lemma is proved.

\medskip

In order to find bounds for the orthogonal polynomials and the kernels, it is
important to know the size of the coefficients which appear in the formulae we
are going to deal with. Frequently, we will restrict ourselves to the case of
measures of compact support and absolutely continuous part almost everywhere
positive.

\proclaim {Lemma 2}. Assume ${\rm supp\ } d\mu = [-1,1]$, $\mu ' > 0$ a.e..
Let $c \in [-1,1]$. Then
$$
   \lim_{n \to \infty} {k_n \over k_n^c} = {1 \over 2},
   \qquad \qquad
   \lim_{n \to \infty} {k_{n-1}^c \over k_{n+1}} = {1 \over 2},
   \qquad \qquad
   \lim_{n \to \infty} {k_n \over k_{n-1}^c} = 1.
$$

\proof The first limit is a consequence of a result of M\'at\'e, Nevai and
Totik (see [3], theorem 11), from which it follows
$$
   \lim_{n \to \infty} {k_n \over k_n^c}
   = \exp (-{1 \over 4\pi}\int_0^{2\pi} \log(\cos t - a)^2 dt).
$$
It is not difficult to see that the integral equals $-4\pi \log 2$, $\forall a
\in [-1,1]$.

The other limits can be obtained from the first one, since from our hypothesis
it follows
$$
   \lim_{n \to \infty} {k_n \over k_{n+1}} = {1 \over 2}
$$
(see [7], pag.~212 or [3], theorem 10). The lemma is proved.

\proclaim {Lemma 3}. Let $d\mu$ be a measure on $\R$ and $n \geq 1$. With the
above notation,
$$
   \int_{\R} P_{n-1}^c(x) (x-c) d\mu(x) =
   - {k_{n-1}^c \over k_n} {P_n(c) \over K_{n-1}(c,c)}
$$

\proof We can write
$$
   P_{n-1}^c(x) (x-c) = \sum_{j=0}^n \alpha_j P_j(x).
$$
By a direct observation of leading coefficients, $\alpha_n = {k_{n-1}^c \over
k_n}$. If $j=1,2,\dots,n-1$, the orthonormality properties of $\{P_n\}$ and
$\{P_n^c\}$ yield
$$ \eqalign{
   \alpha_j &= \int_{\R} P_{n-1}^c(x) (x-c) P_j(x) d\mu(x)
   \cr
   &= \int_{\R} P_{n-1}^c(x) {P_j(x)-P_j(c) \over x-c} (x-c)^2 d\mu(x)
   + \int_{\R} P_{n-1}^c(x) P_j(c) (x-c) d\mu(x)
   \cr
   &= P_j(c) \int_{\R} P_{n-1}^c(x) (x-c) d\mu(x).
   \cr}
$$
For $\alpha_0$ we obtain the same expression, because $P_0(x)$ is a constant:
$$
   \alpha_0 = \int_{\R} P_{n-1}^c(x) (x-c) P_0(x) d\mu(x)
   = P_0(c) \int_{\R} P_{n-1}^c(x) (x-c) d\mu(x).
$$
Therefore
$$ \eqalign{
   P_{n-1}^c(x) (x-c) &= {k_{n-1}^c \over k_n} P_n(x)
   + \sum_{j=0}^{n-1} P_j(c) [\int_{\R} P_{n-1}^c(u) (u-c) d\mu(u)] P_j(x)
   \cr
   &= {k_{n-1}^c \over k_n} P_n(x)
   + K_{n-1}(x,c) \int_{\R} P_{n-1}^c(u) (u-c) d\mu(u).
   \cr}
$$
The lemma follows immediately taking $x=c$ in this equality.

\medskip

We can now obtain an expression for the polynomials orthonormal with respect to
the measure $d\mu + M \delta_c$ in terms of the polynomials $\{P_n\}$ and
$\{P_n^c\}$.

\proclaim {Proposition 4}. Let $d\mu$ be a measure on $\R$, $c \in \R$, $M>0$.
Let $\{Q_n\}_{n\geq 0}$ be the polynomials orthonormal with respect to $d\mu + M
\delta_c$. Then, for each $n \in \N$ there exist two constants $A_n, B_n \in
(0,1)$ such that
$$
   Q_n(x) = A_n P_n(x) + B_n (x-c) P_{n-1}^c(x) \eqno (1)
$$
Furthermore, if ${\rm supp\ }d\mu = [-1,1]$, $\mu ' > 0$ a.e. and $c \in
[-1,1]$, then
$$
   \lim_{n \to \infty} A_n K_{n-1}(c,c) = {1 \over \lambda(c)+M}
$$
and
$$
   \lim_{n \to \infty} B_n = {M \over \lambda(c)+M},
$$
where
$$
   \lambda(c) = \lim_{n \to \infty} {1 \over K_n(c,c)}.
$$

\proof We will find firstly a constant $C_n$ such that $P_n(x) + C_n (x-c)
P_{n-1}^c(x)$ is orthogonal to the polynomials of degree at most $n-1$ with
respect to the measure $d\mu + M \delta_c$. We only need to obtain
$$
   \int_{\R} [P_n(x) + C_n (x-c) P_{n-1}^c(x)] (x-c)^j [d\mu(x) + M \delta_c(x)]
   =0,\quad j=0,1,\dots,n-1. \eqno (2)
$$
Let $j \geq 1$. Then
$$ \eqalign{
   &\int_{\R} [P_n(x) + C_n (x-c) P_{n-1}^c(x)] (x-c)^j [d\mu(x) + M \delta_c(x)]
   \cr
   &= \int_{\R} [P_n(x) + C_n (x-c) P_{n-1}^c(x)] (x-c)^j d\mu(x)
   \cr
   &= \int_{\R} P_n(x) (x-c)^j d\mu(x)
   + C_n \int_{\R} P_{n-1}^c(x) (x-c)^{j-1} d\mu^c(x) = 0.
   \cr}
$$
Therefore, all we have to do is to find a constant $C_n$ for which $(2)$ is
verified with $j=0$. In this case, we can calculate the integral in $(2)$:
$$ \eqalign{
   &\int_{\R} [P_n(x) + C_n (x-c) P_{n-1}^c(x)] [d\mu(x) + M \delta_c(x)]
   \cr
   &= M P_n(c) + \int_{\R} [P_n(x) + C_n (x-c) P_{n-1}^c(x)] d\mu(x)
   \cr
   &= M P_n(c) + C_n \int_{\R} (x-c) P_{n-1}^c(x) d\mu(x)
   = P_n(c) [M - C_n {k_{n-1}^c \over k_n} {1 \over K_{n-1}(c,c)}],
   \cr}
$$
according to lemma $3$. If we take
$$
   C_n = M {k_n \over k_{n-1}^c} K_{n-1}(c,c),
$$
then $P_n(x) + C_n (x-c) P_{n-1}^c(x)$ is orthogonal to every polynomial of
degree at most $n-1$. As $C_n>0$, it is a polynomial of degree $n$ and leading
coefficient positive. Thus, we will obtain the orthonormal polynomial $Q_n$ by
dividing it by its $L^2(d\mu +M\delta_c)$--norm.
$$ \eqalign{
   &\Vert P_n(x) + C_n (x-c) P_{n-1}^c(x) \Vert ^2_{L^2(d\mu +M\delta_c)}
   \cr
   &= M P_n(c)^2 + \int_{\R} [P_n(x) + C_n (x-c) P_{n-1}^c(x)]^2 d\mu(x)
   \cr
   &= M P_n(c)^2 + \int_{\R} P_n(x)^2 d\mu(x)
   + C_n^2 \int_{\R} P_{n-1}^c(x)^2 (x-c)^2 d\mu(x)
   \cr
   &\phantom{= M P_n(c)^2}\;+ 2 C_n \int_{\R} P_n(x) (x-c) P_{n-1}^c(x) d\mu(x)
   \cr
   &= M P_n(c)^2 + 1 + C_n^2 + 2 C_n {k_{n-1}^c \over k_n}.
   \cr}
$$
If we denote
$$
   D_n = [M P_n(c)^2 + 1 + C_n^2 + 2 C_n {k_{n-1}^c \over k_n}]^{1/2},
$$
then we have
$$
   Q_n(x) = {1 \over D_n} P_n(x) + {C_n \over D_n} (x-c) P_{n-1}^c(x),
$$
that is, equation $(1)$ with $A_n = 1/D_n$ and $B_n = C_n/D_n$. From its
definition, it is clear that $D_n > 1$ and $D_n > C_n$, so $A_n, B_n \in (0,1)$.

For the second part, let us assume supp $d\mu = [-1,1]$, $\mu'>0$ a.e. and $c
\in [-1,1]$. From the above definitions for $A_n$, $C_n$ and $D_n$, we have
$$
   {1 \over A_n K_{n-1}(c,c)}
   = [M {P_n(c)^2 \over K_{n-1}(c,c)^2} +
   {1 \over K_{n-1}(c,c)^2} + M^2 ({k_n \over k_{n-1}^c})^2 
   + {2M \over K_{n-1}(c,c)}]^{1/2}.
   \eqno (3)
$$
It can be shown (see [6], theorem $3$, pag.~$26$ and [7], pag.~212 or [3],
theorem 10) that
$$
   \lim_{n \to \infty} {P_n(x)^2 \over K_{n-1}(x,x)} = 0
   \qquad \forall x \in [-1,1].
$$
Since $K_{n-1}(c,c) \geq P_0^2$ this also implies
$$
   \lim_{n \to \infty} {P_n(c)^2 \over K_{n-1}(c,c)^2} = 0.
$$
From this and lemma $2$ we obtain
$$
   \lim_{n \to \infty}{1 \over A_n K_{n-1}(c,c)} = \lambda(c) + M.
$$
Finally,
$$
   \lim_{n \to \infty} B_n
   = \lim_{n \to \infty} A_n M {k_n \over k_{n-1}^c} K_{n-1}(c,c)
   = {M \over \lambda(c) + M}
$$
and the proposition is completely proved.

\medskip

{\noindent \bf Remark.} If $P_n(c)=0$, it is easy to show directly that
$Q_n=P_n$. This is not in contradiction with our proposition, since in this case
it can also be proved that $P_n(x)=(x-c)P_{n-1}^c$ and $A_n + B_n =1$.

\medskip

We can also find some relations which involve the kernels.

\proclaim {Proposition 5}. Let $d\mu$ be a measure on $\R$, $c \in \R$ and
$M>0$. Let $\{L_n\}_{n\geq 0}$ be the kernels relative to $d\mu + M\delta_c$.
Then $\forall n \in \N$
$$
   L_n(x,y) = {1 \over 1 + M K_n(c,c)} K_n(x,y)
   + {M K_n(c,c) \over 1 + M K_n(c,c)} (x-c)(y-c) K_{n-1}^c(x,y)
$$

\proof If $y \in \R$, it is a well-known fact that the kernels $\{K_n^c\}$ verify
$$
   \int_{\R} R_n(x) K_n^c(x,y) (x-c)^2 d\mu(x) = R_n(y)
$$
for every polynomial $R_n$ of degree at most $n$. Actually, this property
characterizes the kernels relative to any measure.

If we write
$$
   (x-c)(y-c)K_{n-1}^c(x,y) = \sum_{j=0}^n \alpha_j(y) P_j(x), \eqno (4)
$$
then it is easy to show that $\forall j \geq 1$
$$
   \alpha_j(y) = (y-c) \int_{\R} K_{n-1}^c(x,y) {P_j(x)-P_j(c) \over x-c}
   (x-c)^2 d\mu(x)
$$
$$
   + (y-c) P_j(c) \int_{\R} K_{n-1}^c(x,y) (x-c) d\mu(x).
$$
By the above property, we obtain
$$
   \alpha_j(y) = P_j(y) - P_j(c)
   + (y-c) P_j(c) \int_{\R} K_{n-1}^c(x,y) (x-c) d\mu(x)
$$
and it is immediate to see that $\alpha_0$ also verifies this formula.

From this formula and $(4)$ it follows
$$ \eqalign{
   (x&-c)(y-c)K_{n-1}^c(x,y)
   \cr
   = K_n(x,y) - K_n(x,c)
   + &(y-c)K_n(x,c)\int_{\R} K_{n-1}^c(u,y) (u-c) d\mu(u)
   \cr}
$$
If we let $x=c$, we obtain
$$
   (y-c) \int_{\R} K_{n-1}^c(x,y) (x-c) d\mu(x) = 1 - {K_n(c,y) \over K_n(c,c)},
$$
and, replacing this equation into the previous one,
$$
   {1 \over 1 + M K_n(c,c)} K_n(x,y) + {M K_n(c,c) \over 1 + M K_n(c,c)}
   (x-c)(y-c) K_{n-1}^c(x,y)
$$
$$
   = K_n(x,y) - {M K_n(c,y) \over 1 + M K_n(c,c)} K_n(x,c).
$$
Therefore, it will be enough to prove that
$$
   \int_{\R} [K_n(x,y) - {M K_n(c,y) \over 1+M K_n(c,c)} K_n(x,c)] R_n(x)
   [d\mu(x) + M \delta_c(x)] = R_n(y)
$$
whenever $R_n$ is a polynomial of degree at most $n$. This an easy consequence
of the fact that the kernels $\{K_n\}$ verify the analogous property with
respect to the measure $d\mu$. The proposition is proved.

\beginsection \S 2. Generalized Jacobi weights with mass points.

Let $\omega$ be a generalized Jacobi weight, that is:
$$
   \omega (x) =h(x)  (1-x)^\alpha(1+x)^\beta\prod_{i=1}^N| x-t_i|^{\gamma_i}\ ,\ 
   x\in [-1,1]
$$
where:

a) $\alpha$,$\beta$,$\gamma_i>-1$, $t_i \in (-1,1)$, $t_i \not= t_j$ $\forall
i \not= j$;

b) $h$ is a positive, continuous function on $[-1,1]$ and
$\omega (h,\delta)\delta^{-1} \in L^1(0,2)$, $\omega (h,\delta)$ being the
modulus of continuity of~$h$.

If we define
$$
   d(x,n) = (1-x+n^{-2})^{-(2\alpha+1)/4}(1+x+n^{-2})^{-(2\beta+1)/4}
   \prod_{i=1}^N (\vert x-t_i \vert + n^{-1})^{-\gamma_i/2},
$$
then the polynomials $\{P_n\}$ orthonormal with respect to the measure $w(x)dx$
on the interval $[-1,1]$ verify the estimate
$$
   \vert P_n(x) \vert \leq C d(x,n) \qquad \forall x \in [-1,1],
   \ \forall n \geq 1, \eqno (5)
$$
where $C$ is a constant independent of $n$ and $x$ (see [2]). In the sequel $C$
will denote a constant independent of $n$ and $x$, but possibly different in
each occurrence.

As to the kernels, it can be shown (see [6], pag.~120 and pag.~4) that
$$
   K_n(x,x) \sim n (1-x+n^{-2})^{-(2\alpha+1)/2} (1+x+n^{-2})^{-(2\beta +1)/2}
   \prod_{i=1}^N (\vert x-t_i \vert+n^{-1})^{-\gamma_i} \eqno (6)
$$
uniformly in $\vert x \vert \leq 1$, $n \geq 1$, where by $f \sim g$ in a domain
$D$ we mean that there exist some positive constants $C_1$ and $C_2$ such that
$C_1 f(y) \leq g(y) \leq C_2 f(y)$ $\forall y \in D$.

Our aim is to prove similar bounds for the polynomials and the kernels relative
to a measure which consists of a generalized Jacobi weight and a finite number
of mass points on the interval $[-1,1]$. So, let $k \in \N$, $a_i \in [-1,1]$ and
$M_i > 0$, $i=1,\dots,k$. We will denote
$$
   d\nu = w(x)dx + \sum_{i=1}^k M_i\delta_{a_i}
$$
on the interval $[-1,1]$. By $\{Q_n\}$ and $\{L_n\}$ we mean, respectively,
their orthonormal polynomials and kernels. Without loss of generality we can
assume $a_i \in \{1,-1,t_1,\dots,t_N\}$, since in the definition of $w$ we can
allow some of the exponents to be $0$. Furthermore, for every $t \in [-1,1]$ we
can speak of its exponent in $w$, referring to the exponent of the factor
$\vert x-t \vert ^\gamma$ in $w$. Obviously, there are only finitely many points
with an exponent different from~$0$.

With this notation, we can deduce some bounds from the results of the previous
section.

\proclaim {Proposition 6}. There exists a constant $C$ such that $\forall n \geq
1$, $\forall x \in [-1,1]$
$$
   \vert Q_n(x) \vert \leq C (1-x+n^{-2})^{-(2\alpha+1)/4}
   (1+x+n^{-2})^{-(2\beta+1)/4} 
   \prod_{i=1}^N (\vert x-t_i \vert+n^{-1})^{-\gamma_i/2},
$$
$$
   \vert L_n(x,x) \vert \leq C (1-x+n^{-2})^{-(2\alpha+1)/2}
   (1+x+n^{-2})^{-(2\beta+1)/2} 
   \prod_{i=1}^N (\vert x-t_i \vert+n^{-1})^{-\gamma_i}.
$$

\proof a) We are going to prove the bound for $Q_n$ by induction on the number
$k$ of mass points. If $k=0$, the measure is a generalized Jacobi weight and we
already know the formula $(5)$. Let $k>0$ and assume the property holds for
$k-1$ mass points.

Let $\{P_n\}$ be the orthonormal polynomials with respect to the measure
$$
   d\mu = w(x)dx + \sum_{i=1}^{k-1} M_i\delta_{a_i}
$$
so that, according to the notation we used in section $1$, $\{P_n^{a_k}\}$ are
the polynomials orthonormal with respect to
$$
   (x-a_k)^2d\mu(x) = (x-a_k)^2w(x)dx
   + \sum_{i=1}^{k-1} (a_i-a_k)^2 M_i\delta_{a_i}.
$$
Since $d\nu = d\mu + M_k\delta_{a_k}$, from proposition $4$ it follows
$$
   Q_n(x) = A_n P_n(x) + B_n (x-a_k) P_{n-1}^{a_k},
$$
with $A_n, B_n \in (0,1)$. Taking into account that both $d\mu$ and
$(x-a_k)^2d\mu(x)$ are generalized Jacobi weights with $k-1$ mass points, they
satisfy the boundedness in the statement. Now, it is easy to see that $Q_n$
satisfies that boundedness.

Therefore, part a) is proved. As to b), proposition $5$ yields
$$
   L_n(x,x) = C_n K_n(x,x) + (1-C_n) (x-a_k)^2 K_{n-1}^{a_k}(x,x)
$$
with $C_n \in (0,1)$. Similar arguments and formula $(6)$ lead to this bound
and the proposition is completely proved.

\medskip

The previous result establishes only upper bounds, which sometimes is not
enough. In some applications (for example, in the study of the convergence
of the Fourier series) it is necessary to estimate more exactly the rate of
growth of $L_n(x,x)$, at least at some points. In the case of a generalized
Jacobi weight, with no point masses, we have even uniform estimates (formula
$(6)$). This estimates cannot hold when the measure has mass points, since there
the kernels $L_n(x,x)$ are bounded (see [6], pag.~4, for example). However, we
can obtain such estimates for each point with no mass.

\proclaim {Proposition 7}. a) Let $t \in (-1,1)$, $t \not= a_i$
$(i=1,\dots,N)$ and let $\gamma$ be its exponent in $w$. Then $L_n(t,t) \sim
n^{1+\gamma}$. \hfil \break
b) Suppose $a_i \not= 1$ $(i=1,\dots,N)$. Then $L_n(1,1) \sim n^{2\alpha
+2}$. \hfil \break
c) Suppose $a_i \not= -1$ $(i=1,\dots,N)$. Then $L_n(-1,-1) \sim n^{2\beta +2}$.

\proof By induction, making use of similar arguments to that of part b) of
proposition~$6$.

\medskip

As an application of the results of section $1$, some bounds for $L_n(x,a_i)$ can
also be obtained. These bounds can be used to study the convergence of the
Fourier series with respect to this type of measure.

\proclaim {Proposition 8}. a) Let $1 \leq i \leq k$ and suppose $a_i \not= \pm
1$. Then there exists a constant $C$ such that $\forall x \in [-1,1]$ and
$\forall n \geq 1$
$$
   \vert L_n(x,a_i) \vert \leq C (1-x+n^{-2})^{-(2\alpha+1)/4}
   (1+x+n^{-2})^{-(2\beta+1)/4} 
   \prod_{t_j \not= a_i} (\vert x-t_j \vert+n^{-1})^{-\gamma_j/2}.
$$
b) If $1$ is a mass point, then there exists a constant $C$ such that $\forall x
\in [-1,1]$ and $\forall n \geq 1$
$$
   \vert L_n(x,1) \vert \leq C (1+x+n^{-2})^{-(2\beta+1)/4} 
   \prod_{i=1}^N (\vert x-t_i \vert+n^{-1})^{-\gamma_i/2}.
$$
c) If $-1$ is a mass point, then there exists a constant $C$ such that $\forall x
\in [-1,1]$ and $\forall n \geq 1$
$$
   \vert L_n(x,-1) \vert \leq C (1-x+n^{-2})^{-(2\alpha+1)/4}
   \prod_{i=1}^N (\vert x-t_i \vert+n^{-1})^{-\gamma_i/2}.
$$

\proof a) Assume $1 \leq i \leq k$ and $a_i \not= \pm 1$. Let $\gamma$ be the
exponent of $\vert x-a_i \vert$ en $w$. If we denote
$$
   d\mu = w(x)dx + \sum_{j=1,j\not= i}^k M_j\delta_{a_j},
$$
then $d\nu = d\mu + M_i\delta_{a_i}$. Let $\{P_n\}$ and $\{K_n\}$ be the
orthonormal polynomials and the kernels relative to $d\mu$ and $k_n$ the
leading coefficient of $P_n$. Analogously, $\{P_n^{a_i}\}$, $\{K_n^{a_i}\}$ and
$\{k_n^{a_i}\}$ with respect to $(x-a_i)^2d\mu$.

If we write
$$
   \Psi_n(x) = (1-x+n^{-2})^{-(2\alpha+1)/4} (1+x+n^{-2})^{-(2\beta +1)/4}
   \prod_{t_j \not= a_i} (\vert x-t_j \vert+n^{-1})^{-\gamma_j/2}
$$
what we have to proof is $\vert L_n(x,a_i) \vert \leq C \Psi_n(x)$. Now, from
proposition $5$ and lemma $1$ we obtain
$$
   L_n(x,a_i) = {k_n \over k_n^{a_i}} {P_n(a_i) \over 1+M_iK_n(a_i,a_i)}
   P_n^{a_i}(x) - {k_{n-1}^{a_i} \over k_{n+1}} {P_{n+1}(a_i) \over
   1+M_iK_n(a_i,a_i)} P_{n-1}^{a_i}(x).
$$
We only need to estimate the right hand side. From proposition $6$ we get
$$ \eqalign{
   \vert P_n(a_i) \vert &\leq C n^{\gamma/2};
   \cr
   \vert P_{n+1}(a_i) \vert &\leq C n^{\gamma/2};
   \cr
   \vert P_n^{a_i}(x) \vert
   &\leq C (\vert x-a_i \vert + n^{-1})^{-(\gamma+2)/2} \Psi_n(x);
   \cr
   \vert P_{n-1}^{a_i}(x) \vert
   &\leq C (\vert x-a_i \vert + n^{-1})^{-(\gamma+2)/2} \Psi_n(x).
   \cr}
$$
Since $a_i$ is not a mass point for $d\mu$, proposition $7$ yields
$$
   K_n(a_i,a_i) \sim n^{1+\gamma}.
$$
Finally, by lemma 2
$$
   \lim_{n \to \infty} {k_n \over k_n^{a_i}} = {1 \over 2}; \qquad
   \lim_{n \to \infty} {k_{n-1}^{a_i} \over k_{n+1}} = {1 \over 2}.
$$
It is now easy to deduce
$$
   \vert L_n(x,a_i) \vert \leq C n^{-1-\gamma/2}(\vert x-a_i \vert +
   n^{-1})^{-1-\gamma/2} \Psi_n(x) \leq C \Psi_n(x).
$$

b) Assume $1$ is a mass point. We define now
$$
   d\mu = w(x)dx + \sum_{i=1,a_i \not= 1}^k M_i\delta_{a_i},
$$
so $d\nu = d\mu + M\delta_1$, $M>0$. If, according to our usual notation,
$\{P_n\}$, $\{K_n\}$ and $\{k_n\}$ refer to $d\mu$ and $\{R_n\}$ are the
orthonormal polynomials relative to the measure $(1-x)d\mu$, $\{r_n\}$ being
their leading coefficients, it is not difficult to show that
$$
   K_n(x,1) = {k_n \over r_n} P_n(1) R_n(x)
$$
(only standard properties of $K_n(x,1)$ are needed). Thus, proposition $5$
leads to
$$
   L_n(x,1) = {k_n \over r_n} {P_n(1) \over 1+M K_n(1,1)} R_n(x).
$$
We proceed now analogously to part a), since $d\mu$ and $(1-x)d\mu$ are
generalized Jacobi weights with masses at points different from $1$. Notice
that, by H\"older's inequaltiy
$$ \eqalign{
   {k_n \over r_n} &= \int_{-1}^1 R_n(x) P_n(x) (1-x) d\mu(x)
   \cr
   &\leq (\int_{-1}^1 R_n(x)^2 (1-x)^2 d\mu(x))^{1/2}
   (\int_{-1}^1 P_n(x)^2 d\mu(x))^{1/2}
   \cr
   &\leq \sqrt 2 (\int_{-1}^1 R_n(x)^2 (1-x)d\mu(x))^{1/2}
   (\int_{-1}^1 P_n(x)^2 d\mu(x))^{1/2} = \sqrt 2.
   \cr}
$$

Part c) is similar to b). Thus, the result is proved.

\bigskip

\centerline {\bf REFERENCES}

\bigskip

\frenchspacing

\art 1;R. Askey and S. Wainger; Mean convergence of expansions in Laguerre and
Hermite series; Amer. J. Math.; 87;1965; 695--708.

\art 2;V. M. Badkov; Convergence in the mean and almost everywhere of Fourier
series in polynomials orthogonal on an interval; Math. USSR Sb.; 24;1974;
223--256.

\art 3;A. M\'at\'e, P. Nevai and V. Totik; Extensions of Szeg\H o's
Theory of Orthogonal Polynomials. II; Constr. Approx.; 3;1987; 51--72.

\art 4;B. Muckenhoupt; Mean convergence of Jacobi series; Proc. Amer. Math.
Soc.; 23;1969; 306--310.

\art 5;B. Muckenhoupt; Mean convergence of Hermite and Laguerre series. II;
Trans. Amer. Math. Soc.; 147;1970; 433--460.

\item {\bf [6]} {\rm P. Nevai, {\it \lq \lq Orthogonal Polynomials"}, Memoirs of
the Amer. Math. Soc., vol.~18, n.~213, Providence, RI, U.S.A., 1979.}

\art 7;E. A. Rahmanov; On the asymptotics of the ratio of orthogonal
polynomials; Math. USSR. Sb.; 32;1977; 199--213.

\item {\bf [8]} {\rm G. Szeg\H o, {\it \lq \lq Orthogonal Polynomials"}, Amer.
Math. Soc. Colloq. Publ., vol.~23, Providence, RI, U.S.A., 1959.}

\bye
